\documentclass[11pt]{article}
\usepackage{amsmath, amssymb, amsthm}
\usepackage[margin=1in]{geometry}
\usepackage{enumitem}

\title{Guidance Approaches for Diffusion-Regularized Utility Optimization}
\author{Discussion Notes}
\date{March 2026}

\newcommand{\xb}{\mathbf{x}}
\newcommand{\xc}{\xb_{\mathrm{c}}}  % current image being optimized
\newcommand{\zb}{\mathbf{z}}
\newcommand{\eb}{\boldsymbol{\epsilon}}
\newcommand{\Ib}{\mathbf{I}}
\newcommand{\R}{\mathbb{R}}
\newcommand{\E}{\mathbb{E}}
\newcommand{\alphabar}{\bar{\alpha}}

\begin{document}
\maketitle

\section{Setup and Notation}

We are optimizing an image $\xc \in \R^d$ to maximize a GP-based utility function $U(\xc)$
while keeping $\xc$ on (or near) the natural image manifold $\mathcal{M}$.

At each optimization step, we update $\xc$ by following the gradient of a combined objective.
The utility gradient $\nabla_{\xc} U(\xc)$ pushes $\xc$ toward high-utility regions, but tends to
drive it off $\mathcal{M}$ (e.g., by amplifying pixel norms). We want to add a ``naturalness'' term
that counteracts this drift.

We have a trained DDPM with noise predictor $\eb_\theta(\xb_t, t)$, trained on natural images
corrupted at noise levels $t = 1, \ldots, T$. The noise schedule defines $\alphabar_t = \prod_{s=1}^t (1-\beta_s)$.

The key identity relating noise prediction to the score (gradient of log-density):
\begin{equation}
    \nabla_{\xb_t} \log p_t(\xb_t) \approx s_\theta(\xb_t, t)
    = -\frac{\eb_\theta(\xb_t, t)}{\sqrt{1 - \alphabar_t}}
    \label{eq:score_from_eps}
\end{equation}

And the Tweedie estimate of the clean image given a noisy observation:
\begin{equation}
    \hat{\xb}_0(\xb_t, t)
    = \frac{\xb_t - \sqrt{1 - \alphabar_t}\, \eb_\theta(\xb_t, t)}{\sqrt{\alphabar_t}}
    \label{eq:tweedie}
\end{equation}

Below we discuss three approaches to incorporating the diffusion model into the optimization,
with their respective problems.


% ============================================================================
\section{Approach A: L2 Penalty Toward Tweedie Estimate (Current Implementation)}
% ============================================================================

\subsection{Method}

At each optimization step:
\begin{enumerate}[nosep]
    \item Treat $\xc$ as a clean image $\xb_0$.
    \item Corrupt it at a fixed noise level $t^*$ (e.g., $t^* = 50$):
    \begin{equation}
        \xb_t = \sqrt{\alphabar_{t^*}}\, \xc + \sqrt{1 - \alphabar_{t^*}}\, \eb,
        \qquad \eb \sim \mathcal{N}(\mathbf{0}, \Ib)
        \label{eq:corrupt}
    \end{equation}
    \item Feed $\xb_t$ to the UNet to predict noise $\hat{\eb} = \eb_\theta(\xb_t, t^*)$.
    \item Compute the Tweedie estimate via Eq.~\eqref{eq:tweedie}:
    \begin{equation}
        \hat{\xb}_0 = \frac{\xb_t - \sqrt{1-\alphabar_{t^*}}\,\hat{\eb}}{\sqrt{\alphabar_{t^*}}}
    \end{equation}
    \item \textbf{Detach} $\hat{\xb}_0$ from the computation graph (no backprop through UNet).
    \item Add an L2 penalty pulling $\xc$ toward $\hat{\xb}_0$:
    \begin{equation}
        \mathcal{L} = -U(\xc) + \frac{\lambda}{2} \|\xc - \hat{\xb}_0\|^2
        \label{eq:loss_l2}
    \end{equation}
\end{enumerate}

The gradient of the regularization term w.r.t.\ $\xc$ is:
\begin{equation}
    \nabla_{\xc} \left[\frac{\lambda}{2}\|\xc - \hat{\xb}_0\|^2\right]
    = \lambda\,(\xc - \hat{\xb}_0)
\end{equation}

The SGD update becomes:
\begin{equation}
    \xc \leftarrow \xc + \eta\,\nabla_{\xc} U(\xc)
    + \eta\,\lambda\,(\hat{\xb}_0 - \xc)
    \label{eq:update_l2}
\end{equation}

\subsection{Connection to Score Distillation}

Substituting the definitions of $\xb_t$ and $\hat{\xb}_0$, the regularization gradient simplifies to:
\begin{equation}
    \xc - \hat{\xb}_0
    = \frac{\sqrt{1-\alphabar_{t^*}}}{\sqrt{\alphabar_{t^*}}}\,(\hat{\eb} - \eb)
    \label{eq:sds_connection}
\end{equation}

This is proportional to $(\hat{\eb} - \eb)$, the difference between the predicted and actual noise.
This is the same quantity used in Score Distillation Sampling (SDS), though with a different
weighting factor. SDS uses $w(t)\sqrt{\alphabar_t}\,(\hat{\eb} - \eb)$; our L2 formulation
uses $\sqrt{(1-\alphabar_t)/\alphabar_t}\,(\hat{\eb} - \eb)$.

\subsection{Problems}

\begin{enumerate}
    \item \textbf{No signal near the natural manifold.}
    When $\xc$ is a natural image, $\hat{\xb}_0 \approx \xc$ (the UNet correctly recovers
    the clean image from its noisy version). The residual $\xc - \hat{\xb}_0$ is dominated
    by the UNet's reconstruction error, which is effectively random noise with no systematic
    direction. The regularization provides no useful gradient near $\mathcal{M}$.

    This is by design: the penalty is \emph{adaptive} --- it is small when $\xc$ is natural
    and grows as $\xc$ drifts away. But it means the method has no preventive force; it only
    reacts after the image has already become unnatural.

    \item \textbf{Norm amplification not prevented.}
    The Tweedie estimate $\hat{\xb}_0$ of an amplified image is itself amplified. The UNet
    denoises ``amplified + noise'' back to ``amplified clean.'' The L2 penalty preserves
    structure (spatial correlations) but does not penalize the overall amplitude. The direction
    $(\hat{\xb}_0 - \xc)$ corrects local texture but not global norm.

    \item \textbf{Fixed noise level $t^*$.}
    The choice of $t^*$ determines what scale of features the regularization ``sees.''
    At $t^* = 50$ ($\alphabar \approx 0.986$, noise $\approx 12\%$ of signal), the UNet
    captures fine-to-medium structure. Large-scale distortions (e.g., global brightening)
    may be invisible at this noise level. There is no principled way to choose $t^*$.

    \item \textbf{Single noise sample.}
    Each step uses one draw of $\eb$. The gradient $(\hat{\eb} - \eb)$ is a single-sample
    estimate, introducing variance. Using a fixed seed (as in the current implementation)
    removes variance but introduces bias: the regularization always ``sees'' the same
    noise realization.
\end{enumerate}


% ============================================================================
\section{Approach B: Direct Score as Gradient Term}
% ============================================================================

\subsection{Method}

Use the score function $s_\theta(\xb_t, t)$ directly as a gradient contribution.
At each optimization step:
\begin{enumerate}[nosep]
    \item Corrupt $\xc$ at noise level $t^*$:
    $\xb_t = \sqrt{\alphabar_{t^*}}\, \xc + \sqrt{1 - \alphabar_{t^*}}\, \eb$
    \item Predict noise: $\hat{\eb} = \eb_\theta(\xb_t, t^*)$
    \item Compute the score via Eq.~\eqref{eq:score_from_eps}:
    \begin{equation}
        s_\theta(\xb_t, t^*) = -\frac{\hat{\eb}}{\sqrt{1-\alphabar_{t^*}}}
    \end{equation}
    \item Use the score as the naturalness gradient, scaled to $\xc$-space by the
    chain rule $\partial \xb_t / \partial \xc = \sqrt{\alphabar_{t^*}}$:
    \begin{equation}
        \nabla_{\xc} \log p_{t^*}(\xb_t) \approx \sqrt{\alphabar_{t^*}}\, s_\theta(\xb_t, t^*)
        = -\frac{\sqrt{\alphabar_{t^*}}}{\sqrt{1-\alphabar_{t^*}}}\, \hat{\eb}
    \end{equation}
\end{enumerate}

The update becomes:
\begin{equation}
    \xc \leftarrow \xc + \eta\,\nabla_{\xc} U(\xc)
    + \eta\, w \,\frac{\sqrt{\alphabar_{t^*}}}{\sqrt{1-\alphabar_{t^*}}}\,
    (-\hat{\eb})
    \label{eq:update_score}
\end{equation}

where $w$ is the guidance weight.

\subsection{Comparison with Approach A}

The key difference: Approach~A uses $(\hat{\eb} - \eb)$ (prediction minus actual noise);
Approach~B uses $\hat{\eb}$ alone (the full prediction). Approach~B includes the actual
noise direction as part of its gradient, which adds variance but provides signal even
when $\xc$ is on the natural manifold.

When $\xc$ is natural and the UNet predicts perfectly ($\hat{\eb} \approx \eb$):
\begin{itemize}[nosep]
    \item Approach A gradient: $\propto (\hat{\eb} - \eb) \approx 0$ (no signal).
    \item Approach B gradient: $\propto \hat{\eb} \approx \eb$ (nonzero, but random direction).
\end{itemize}

Neither provides a useful deterministic direction for natural images. However, in expectation
over $\eb$, the Approach~B gradient $\E_\eb[-\hat{\eb}]$ is nonzero when $\xc$ is not at a mode
of $p_{t^*}$, because the score has nonzero expectation there.

\subsection{Problems}

\begin{enumerate}
    \item \textbf{Same noise requirement as Approach A.}
    The UNet still requires noise-corrupted input. We must add artificial noise
    at level $t^*$ and accept that we get $\nabla \log p_{t^*}$, not $\nabla \log p_0$.

    \item \textbf{High variance from the noise term.}
    The gradient includes $\hat{\eb}$, which contains the actual noise $\eb$ as a large
    component. This makes each gradient step noisy. In Approach~A, the noise cancels
    ($\hat{\eb} - \eb$), reducing variance. In Approach~B it does not.

    \item \textbf{Score of $p_{t^*}$, not $p_0$.}
    The score $s_\theta(\xb_t, t^*)$ estimates $\nabla \log p_{t^*}(\xb_t)$, which is
    the gradient of the log-density of the \emph{noise-smoothed} distribution at level $t^*$.
    This is a blurred version of $\nabla \log p_0(\xc)$. Fine-scale naturalness features
    (edges, textures) are washed out at moderate $t^*$.

    \item \textbf{Scale mismatch with utility gradient.}
    The score magnitude depends on $t^*$ through the $1/\sqrt{1-\alphabar_{t^*}}$ factor.
    At small $t^*$ this factor is large (e.g., $\sim 100$ at $t=1$), at large $t^*$
    it is $O(1)$. Balancing the score against the utility gradient requires careful
    tuning of both $w$ and $t^*$.
\end{enumerate}


% ============================================================================
\section{Approach C: Score Without Noise (Direct Evaluation at $\xc$)}
% ============================================================================

\subsection{Method}

Feed $\xc$ directly to the UNet without adding noise, using a small timestep $t = 1$:
\begin{enumerate}[nosep]
    \item Pass $(\xc, t{=}1)$ to the UNet: $\hat{\eb} = \eb_\theta(\xc, 1)$.
    \item Compute the score: $s_\theta(\xc, 1) = -\hat{\eb} / \sqrt{1-\alphabar_1}$.
    \item Use as gradient:
    \begin{equation}
        \xc \leftarrow \xc + \eta\,\nabla_{\xc} U(\xc)
        + \eta\, w\, s_\theta(\xc, 1)
    \end{equation}
\end{enumerate}

The appeal: no artificial noise, no Tweedie formula. The score would directly approximate
$\nabla_{\xc} \log p_0(\xc)$ since $p_1 \approx p_0$.

\subsection{Problems}

\begin{enumerate}
    \item \textbf{Out-of-distribution input.}
    At $t=1$, $\alphabar_1 \approx 0.9999$, so the UNet expects input of the form
    ``natural image $+$ noise with $\sigma \approx 0.01$.'' During training, every
    $({\xb_t}, t{=}1)$ pair the network saw was a nearly-clean natural image with
    imperceptible perturbation.

    If $\xc$ has been distorted by utility optimization (amplified norms, broken
    correlations), it looks nothing like what the UNet saw at $t=1$ during training.
    The network must extrapolate, and $\hat{\eb}$ is unreliable.

    \item \textbf{Scale mismatch between distortion and noise level.}
    At $t=1$, the expected noise magnitude is $\sqrt{1 - \alphabar_1} \approx 0.01$.
    The UNet is trained to detect perturbations at this tiny scale. Utility-driven
    distortions operate at scale $O(1)$ or larger (e.g., $3\times$ amplification).
    The UNet cannot ``see'' these large-scale distortions as something to correct
    at $t=1$, because they are orders of magnitude larger than the noise budget.

    The unnaturalness is invisible to the network at this noise level.

    \item \textbf{Amplification by $1/\sqrt{1-\alphabar_1}$.}
    Dividing by $\sqrt{1 - \alphabar_1} \approx 0.01$ amplifies $\hat{\eb}$ by
    $\sim 100\times$. Since $\hat{\eb}$ is unreliable (problem 1), this produces
    a large, meaningless gradient vector.

    \item \textbf{Theoretical vs.\ practical.}
    In theory, $\nabla \log p_1(\xc) \approx \nabla \log p_0(\xc)$ for any $\xc$,
    including unnatural ones. This \emph{would} point toward naturalness. But the
    UNet's approximation $s_\theta(\xc, 1)$ is only reliable for inputs near the
    training distribution. Unnatural images at $t=1$ are outside this distribution,
    so the approximation breaks down precisely where it would be most needed.
\end{enumerate}


% ============================================================================
\section{Approach D: Full Guided Reverse Diffusion (Reference Method)}
% ============================================================================

\subsection{Method}

Instead of optimizing $\xc$ by gradient descent, generate the image using the full
reverse diffusion process with classifier guidance, as described in the companion document
(\texttt{diffusion\_models\_introduction.tex}, Section~6.3).

\begin{enumerate}[nosep]
    \item Start from noise: $\xb_T \sim \mathcal{N}(\mathbf{0}, \Ib)$.
    \item For $t = T, T{-}1, \ldots, 1$:
    \begin{enumerate}[nosep]
        \item Predict noise: $\hat{\eb} = \eb_\theta(\xb_t, t)$.
        \item Compute the Tweedie estimate: $\hat{\xb}_0$ via Eq.~\eqref{eq:tweedie}.
        \item Compute utility gradient on the Tweedie estimate:
              $\mathbf{g} = \nabla_{\xb_t} U(\hat{\xb}_0)$
              (requires backprop through the Tweedie formula into the UNet).
        \item Modify noise prediction:
              $\hat{\eb}_{\mathrm{guided}} = \hat{\eb} - w\sqrt{1-\alphabar_t}\,\mathbf{g}$.
        \item Reverse step:
              $\xb_{t-1} = \mu_\theta(\xb_t, t;\, \hat{\eb}_{\mathrm{guided}}) + \sigma_t \zb$,
              $\zb \sim \mathcal{N}(\mathbf{0}, \Ib)$.
    \end{enumerate}
    \item Return $\xb_0$.
\end{enumerate}

The output is approximately a sample from $p(\xb \mid \text{high } U)$.

\subsection{Advantages Over Approaches A--C}

\begin{itemize}[nosep]
    \item The score is always evaluated at the correct noise level: at step $t$,
          $\xb_t$ is genuinely at noise level $t$ (it came from the reverse process),
          so the UNet operates within its training distribution.
    \item The process naturally traverses from coarse ($t$ large) to fine ($t$ small)
          structure, capturing naturalness at all scales.
    \item No fixed $t^*$ to choose --- all noise levels are used.
    \item Output is a proper sample from (an approximation of) the guided distribution,
          not a point estimate from gradient descent.
\end{itemize}

\subsection{Problems}

\begin{enumerate}
    \item \textbf{Requires backprop through UNet.}
    The utility gradient $\mathbf{g} = \nabla_{\xb_t} U(\hat{\xb}_0)$ requires
    differentiating through the Tweedie formula and the UNet forward pass. This means
    the UNet must run with gradients enabled (not just \texttt{torch.no\_grad}).
    Memory and compute cost per step increase significantly.

    \item \textbf{$T$ sequential UNet evaluations.}
    Full reverse diffusion requires $T$ steps (or $\sim 50$--$200$ with accelerated
    samplers like DDIM). Each step involves a UNet forward pass plus a utility
    evaluation (which itself requires a GP forward pass). Total cost is $O(T)$ times
    the cost of a single optimization step in Approaches A--C.

    \item \textbf{Utility on noisy intermediates.}
    $U(\hat{\xb}_0)$ evaluates the utility on the Tweedie estimate, which is a rough
    approximation of the final image (especially at large $t$). Early guidance steps
    ($t$ near $T$) operate on very noisy estimates and may provide misleading gradients
    for the utility.

    \item \textbf{Stochastic output.}
    Each run produces a different image (different $\xb_T$, different noise in reverse
    steps). To find the best image, multiple runs are needed. This is a feature
    (diversity) but also a cost.

    \item \textbf{Starting point not controllable.}
    Unlike Approaches A--C, which start from a chosen image (e.g., smoothed target),
    Approach~D starts from pure noise. The generated image will be ``natural with high
    utility'' but may not resemble any specific target image.
\end{enumerate}


% ============================================================================
\section{Summary}
% ============================================================================

\begin{center}
\begin{tabular}{l c c c c}
\hline
& \textbf{A: L2 Tweedie} & \textbf{B: Direct Score} & \textbf{C: No Noise} & \textbf{D: Guided Sampling} \\
\hline
Noise required          & Yes ($t^*$)   & Yes ($t^*$)   & No            & Inherent \\
Backprop through UNet   & No            & No            & No            & Yes \\
Score of $p_0$ or $p_t$ & Neither (L2)  & $p_{t^*}$     & $p_1$ (unreliable) & All $p_t$ \\
Signal near $\mathcal{M}$ & Weak        & Noisy         & Unreliable    & Strong \\
Prevents norm growth    & No            & Partially     & Unreliable    & Yes \\
Cost per step           & 1 UNet fwd    & 1 UNet fwd    & 1 UNet fwd    & $T$ UNet fwd+bwd \\
Starting point          & Chosen        & Chosen        & Chosen        & Random noise \\
\hline
\end{tabular}
\end{center}

Approach~A is the simplest and cheapest. Its main weakness is the lack of signal near
$\mathcal{M}$ and the inability to prevent norm amplification.

Approach~D is the most principled but the most expensive. It requires backpropagation
through the UNet and $T$ sequential steps per generated image.

Approaches~B and C are intermediate: they attempt to use the score more directly but
face either high variance (B) or out-of-distribution issues (C).

\end{document}